% ============================================================
%  1. INTRODUCCIÓN   (guía, 2.1)
%  Propósito: presentar el tema, el contexto, la importancia del
%  trabajo y la solución que será desarrollada.
%  Debe responder: ¿sobre qué trata el proyecto?, ¿en qué contexto
%  se desarrolla?, ¿por qué es relevante para Ciencia de Datos?,
%  ¿qué encontrará el lector en la monografía?
%  Recomendación: comprensible para un lector que aún no conoce
%  los detalles técnicos.
% ============================================================
\chapter{Introducción}
\label{cap:introduccion}

El transporte público del eje metropolitano de Cochabamba opera principalmente bajo la
modalidad de paratránsito. Trufi App constituye una de las pocas fuentes públicas de
información sobre sus rutas, mantenida mediante el trabajo voluntario de Trufi Association.
La aplicación registra cuántas consultas de ruta se originan en cada zona, pero no cuántas
cabría esperar según sus características territoriales; en consecuencia, un conteo bajo no
permite distinguir entre una zona poco poblada y una zona donde la aplicación no resulta
útil.

La presente monografía aborda ese problema desde la Ciencia de Datos geoespacial. Su
objetivo es estimar el número esperado de consultas de ruta por celda H3 y compararlo con
las consultas efectivamente registradas, a fin de identificar zonas con demanda de
información no resuelta. El proyecto se desarrolla bajo la metodología CRISP-DM, con
validación espacial por bloques y una regla de selección de modelo declarada antes de
conocer los resultados.

El documento se estructura de forma secuencial. Los capítulos~\ref{cap:problema}
a~\ref{cap:objetivos} identifican el problema, justifican su estudio, delimitan el alcance y
establecen los objetivos. El capítulo~\ref{cap:marco-teorico} presenta los fundamentos
teóricos. El capítulo~\ref{cap:desarrollo} documenta la aplicación de las seis fases de
CRISP-DM. El capítulo~\ref{cap:conclusiones} expone las conclusiones y recomendaciones. El
proyecto no contempla despliegue operativo: se limita a un prototipo verificable localmente,
un archivo de predicciones por celda y una lista priorizada de veinte celdas para revisión
en terreno.
